% ── Zaměstnavatelská organizovanost – mapa + vývoj (stav_zamestnavatele_hustota.py) ──

Vedle odborové organizovanosti a~pokrytí \ac{KS} představuje zaměstnavatelská organizovanost třetí klíčový ukazatel organizovanosti sociálních partnerů. V~databázi \ac{OECD}~\ac{ICTWSS} odpovídá ukazateli \emph{employer organisation density} (ED) a~udává podíl zaměstnanců pracujících v~podnicích, které jsou členy zaměstnavatelského svazu.~\cite{oecd_aias_ictwss_ED_pct}
Vysoká zaměstnavatelská organizovanost je předpokladem sektorového \ac{KV} --- bez reprezentativního protějšku na straně zaměstnavatelů nemohou odbory uzavírat odvětvové \ac{KS} s~plošným dopadem.

Obr.~\ref{fig:stav_zamestnavatele_hustota_mapa} ukazuje geografické rozložení zaměstnavatelské organizovanosti v~\ac{EU}.
Rakousko, Nizozemsko a~skandinávské státy vykazují hodnoty přesahující \SI{60}{\percent}, zatímco středoevropské a~baltské země se pohybují výrazně níže.
ČR se řadí mezi státy se střední zaměstnavatelskou organizovaností, přičemž hodnota nedosahuje úrovně potřebné pro plně funkční sektorové vyjednávání.

% Editable figure definition for stav_zamestnavatele_hustota_mapa
% Edit freely — Python will NOT overwrite this file.
% To regenerate defaults, delete this file and re-run the script.
\def\strStavZamestnavateleHustotaMapaCaption{\centering Hustota zaměstnavatelských organizací, EU mapa, 2024. Zdroj dat: OECD\,ICTWSS~\cite{oecd_aias_ictwss_ED_pct}.}%
%
\begin{figure}[htbp]
  \centering
  \resizebox{0.92\linewidth}{!}{% Editable figure definition for stav_zamestnavatele_hustota_mapa
% Edit freely — Python will NOT overwrite this file.
% To regenerate defaults, delete this file and re-run the script.
\def\strStavZamestnavateleHustotaMapaCaption{\centering Hustota zaměstnavatelských organizací, EU mapa, 2024. Zdroj dat: OECD\,ICTWSS~\cite{oecd_aias_ictwss_ED_pct}.}%
%
\begin{figure}[htbp]
  \centering
  \resizebox{0.92\linewidth}{!}{% Editable figure definition for stav_zamestnavatele_hustota_mapa
% Edit freely — Python will NOT overwrite this file.
% To regenerate defaults, delete this file and re-run the script.
\def\strStavZamestnavateleHustotaMapaCaption{\centering Hustota zaměstnavatelských organizací, EU mapa, 2024. Zdroj dat: OECD\,ICTWSS~\cite{oecd_aias_ictwss_ED_pct}.}%
%
\begin{figure}[htbp]
  \centering
  \resizebox{0.92\linewidth}{!}{\input{../python/figures/stav_zamestnavatele_hustota_mapa.pgf}}
  \caption{\strStavZamestnavateleHustotaMapaCaption}
  \label{fig:stav_zamestnavatele_hustota_mapa}
\end{figure}
}
  \caption{\strStavZamestnavateleHustotaMapaCaption}
  \label{fig:stav_zamestnavatele_hustota_mapa}
\end{figure}
}
  \caption{\strStavZamestnavateleHustotaMapaCaption}
  \label{fig:stav_zamestnavatele_hustota_mapa}
\end{figure}


Časový vývoj (obr.~\ref{fig:stav_zamestnavatele_hustota_vyvoj}) potvrzuje, že organizovanost zaměstnavatelů se v~jednotlivých zemích vyvíjí odlišně.
Rakousko a~Dánsko si udržují stabilně vysokou zaměstnavatelskou organizovanost, zatímco Polsko a~Slovensko zaznamenaly pokles.
Klíčový poznatek pro ČR spočívá v~tom, že pokles zaměstnavatelské organizovanosti probíhá paralelně s~poklesem odborové organizovanosti --- oba pilíře sociálního dialogu tedy slábnou současně, což ohrožuje funkčnost celého systému \ac{KV}.
Posílení zaměstnavatelské organizovanosti --- např.\ prostřednictvím povinného členství v~hospodářských komorách po vzoru Rakouska --- je proto stejně důležité jako podpora odborů.

% Editable figure definition for stav_zamestnavatele_hustota_vyvoj
% Edit freely --- Python will NOT overwrite this file.
% To regenerate defaults, delete this file and re-run the script.
\def\strStavZamestnavateleHustotaVyvojTitle{Zaměstnavatelská organizovanost}%
\def\strStavZamestnavateleHustotaVyvojYlabel{zaměstnavatelská organizovanost [\%]}%
% --- Per-label y-nudge knobs (override with \renewcommand)
% Example: \renewcommand\NudgeStavZamestnavateleHustotaVyvojCz{-3pt}  % shift label up by 3pt
\providecommand\NudgeStavZamestnavateleHustotaVyvojCz{0pt}%
\providecommand\NudgeStavZamestnavateleHustotaVyvojAt{0pt}%
\providecommand\NudgeStavZamestnavateleHustotaVyvojDe{0pt}%
\providecommand\NudgeStavZamestnavateleHustotaVyvojDk{0pt}%
\providecommand\NudgeStavZamestnavateleHustotaVyvojPl{0pt}%
\providecommand\NudgeStavZamestnavateleHustotaVyvojSk{0pt}%
\def\strStavZamestnavateleHustotaVyvojCaption{Vývoj zaměstnavatelské organizovanosti, vybrané země EU, 2000--2024. Zdroj dat: \acs{OECD}~\acs{ICTWSS}~\cite{oecd_aias_ictwss_ED_pct}.}%
%
\begin{figure}[H]
  \centering
  \resizebox{\linewidth}{!}{% Editable figure definition for stav_zamestnavatele_hustota_vyvoj
% Edit freely --- Python will NOT overwrite this file.
% To regenerate defaults, delete this file and re-run the script.
\def\strStavZamestnavateleHustotaVyvojTitle{Zaměstnavatelská organizovanost}%
\def\strStavZamestnavateleHustotaVyvojYlabel{zaměstnavatelská organizovanost [\%]}%
% --- Per-label y-nudge knobs (override with \renewcommand)
% Example: \renewcommand\NudgeStavZamestnavateleHustotaVyvojCz{-3pt}  % shift label up by 3pt
\providecommand\NudgeStavZamestnavateleHustotaVyvojCz{0pt}%
\providecommand\NudgeStavZamestnavateleHustotaVyvojAt{0pt}%
\providecommand\NudgeStavZamestnavateleHustotaVyvojDe{0pt}%
\providecommand\NudgeStavZamestnavateleHustotaVyvojDk{0pt}%
\providecommand\NudgeStavZamestnavateleHustotaVyvojPl{0pt}%
\providecommand\NudgeStavZamestnavateleHustotaVyvojSk{0pt}%
\def\strStavZamestnavateleHustotaVyvojCaption{Vývoj zaměstnavatelské organizovanosti, vybrané země EU, 2000--2024. Zdroj dat: \acs{OECD}~\acs{ICTWSS}~\cite{oecd_aias_ictwss_ED_pct}.}%
%
\begin{figure}[H]
  \centering
  \resizebox{\linewidth}{!}{% Editable figure definition for stav_zamestnavatele_hustota_vyvoj
% Edit freely --- Python will NOT overwrite this file.
% To regenerate defaults, delete this file and re-run the script.
\def\strStavZamestnavateleHustotaVyvojTitle{Zaměstnavatelská organizovanost}%
\def\strStavZamestnavateleHustotaVyvojYlabel{zaměstnavatelská organizovanost [\%]}%
% --- Per-label y-nudge knobs (override with \renewcommand)
% Example: \renewcommand\NudgeStavZamestnavateleHustotaVyvojCz{-3pt}  % shift label up by 3pt
\providecommand\NudgeStavZamestnavateleHustotaVyvojCz{0pt}%
\providecommand\NudgeStavZamestnavateleHustotaVyvojAt{0pt}%
\providecommand\NudgeStavZamestnavateleHustotaVyvojDe{0pt}%
\providecommand\NudgeStavZamestnavateleHustotaVyvojDk{0pt}%
\providecommand\NudgeStavZamestnavateleHustotaVyvojPl{0pt}%
\providecommand\NudgeStavZamestnavateleHustotaVyvojSk{0pt}%
\def\strStavZamestnavateleHustotaVyvojCaption{Vývoj zaměstnavatelské organizovanosti, vybrané země EU, 2000--2024. Zdroj dat: \acs{OECD}~\acs{ICTWSS}~\cite{oecd_aias_ictwss_ED_pct}.}%
%
\begin{figure}[H]
  \centering
  \resizebox{\linewidth}{!}{\input{../python/figures/stav_zamestnavatele_hustota_vyvoj.pgf}}
  \caption{\strStavZamestnavateleHustotaVyvojCaption}
  \label{fig:stav_zamestnavatele_hustota_vyvoj}
\end{figure}
}
  \caption{\strStavZamestnavateleHustotaVyvojCaption}
  \label{fig:stav_zamestnavatele_hustota_vyvoj}
\end{figure}
}
  \caption{\strStavZamestnavateleHustotaVyvojCaption}
  \label{fig:stav_zamestnavatele_hustota_vyvoj}
\end{figure}

